\chapter{2009年上海卷（秋文）}

\section{填空题}

\begin{problem}
函数 \(f(x) = x^3 + 1\) 的反函数 \(f^{-1}(x) =\)\fillinblank{}.

\end{problem}

\begin{problem}
已知集合 \(A = \{x \mid x \leqslant 1\}\)，\(B = \{x \mid x \geqslant a\}\)，且 \(A \cup B = \R\)，则实数 \(a\) 的取值范围是 \fillinblank{}.

\end{problem}

\begin{problem}
若行列式 \(\begin{vmatrix} 4 & 5 & x \\ 1 & x & 3 \\ 7 & 8 & 9 \end{vmatrix}\) 中，元素 4 的代数余子式大于 0，则 \(x\) 满足的条件是 \fillinblank{}.

\end{problem}

\begin{problem}
某算法的程序框如图所示，则输出量 \(y\) 与输入量 \(x\) 满足的关系式是 \fillinblank{}.

\bitmapfigure[max width=0.30\linewidth,max totalheight=4.2cm]{img/2009/shanghai_liberal/q4_fig1.png}

\end{problem}

\begin{problem}
如图，若正四棱柱 \(ABCD - A_1B_1C_1D_1\) 的底面边长为 2，高为 4，则异面直线 \(BD_1\) 与 \(AD\) 所成角的大小是 \fillinblank{}. （结果用反三角函数表示）

\bitmapfigure[max width=0.35\linewidth,max totalheight=5.0cm]{img/2009/shanghai_liberal/q5_fig1.png}

\end{problem}

\begin{problem}
若球 \(O_1\)、\(O_2\) 表面积之比 \(\frac{S_1}{S_2} = 4\)，则它们的半径之比 \(\frac{R_1}{R_2} =\)\fillinblank{}.

\end{problem}

\begin{problem}
已知实数 \(x\)、\(y\) 满足 \(\left\{ \begin{array}{l} y \leqslant 2x, \\ y \geqslant -2x, \\ x \leqslant 3, \end{array} \right.\) 则目标函数 \(z = x - 2y\) 的最小值是 \fillinblank{}.

\end{problem}

\begin{problem}
若等腰直角三角形的直角边长为 2，则以一直角边所在的直线为轴旋转一周所成的几何体体积是 \fillinblank{}.

\end{problem}

\begin{problem}
过点 \(A(1, 0)\) 作倾斜角为 \(\frac{\pi}{4}\) 的直线，与抛物线 \(y^2 = 2x\) 交于 \(M\)、\(N\) 两点，则 \(|MN| =\)\fillinblank{}.

\end{problem}

\begin{problem}
函数 \(y = 2\cos^2 x + \sin 2x\) 的最小值是 \fillinblank{}.

\end{problem}

\begin{problem}
若某学校要从 5 名男生和 2 名女生中选出 3 人作为上海世博会的志愿者，则选出的志愿者中男女生均不少于 1 名的概率是 \fillinblank{}. （结果用最简分数表示）

\end{problem}

\begin{problem}
已知 \(F_1\)、\(F_2\) 是椭圆 \(C: \frac{x^2}{a^2} + \frac{y^2}{b^2} = 1\)（\(a > b > 0\)）的两个焦点，\(P\) 为椭圆 \(C\) 上一点，且 \(\overrightarrow{PF_1} \perp \overrightarrow{PF_2}\). 若 \(\triangle PF_1F_2\) 的面积为 9，则 \(b =\)\fillinblank{}.

\end{problem}

\begin{problem}
已知函数 \(f(x) = \sin x + \tan x\). 项数为 27 的等差数列 \(\{a_n\}\) 满足 \(a_n \in \left( -\frac{\pi}{2}, \frac{\pi}{2} \right)\)，且公差 \(d \neq 0\). 若 \(f(a_1) + f(a_2) + \cdots + f(a_{27}) = 0\)，则当 \(k =\)\fillinblank{} 时，\(f(a_k) = 0\).

\end{problem}

\begin{problem}
某地街道呈现东一西、南一北向的网格状，相邻街距都为 1，两街道相交的点称为格点. 若以互相垂直的两条街道为轴建立直角坐标系，现有下述格点 \((-2, 2)\)，\((3, 1)\)，\((3, 4)\)，\((-2, 3)\)，\((4, 5)\) 为报刊零售点. 请确定一个格点（除零售点外）为发行站，使 5 个零售点沿街道到发行站之间路程的和最短.

\end{problem}

\setcounter{probnum}{18}
\begin{problem}
已知复数 \(z = a + b\i\)（\(a, b \in \R^+\)）（\(i\) 是虚数单位）是方程 \(x^2 - 4x + 5 = 0\) 的根. 复数 \(w = u + 3\i\)（\(u \in \R\)）满足 \(|w - z| < 2\sqrt{5}\)，求 \(u\) 的取值范围 \fillinblank{}.

\end{problem}

\setcounter{probnum}{14}
\section{单选题}

\begin{problem}
已知直线 \(l_1: (k - 3)x + (4 - k)y + 1 = 0\) 与 \(l_2: 2(k - 3)x - 2y + 3 = 0\) 平行，则 \(k\) 的值是（\quad）

\choices
    {1 或 3}
    {1 或 5}
    {3 或 5}
    {1 或 2}

\end{problem}

\begin{problem}
如图，已知三棱锥的底面是直角三角形，直角边长分别为 3 和 4，过直角顶点的侧棱长为 4，且垂直于底面，该三棱锥的主视图是（\quad）

\bitmapfigure[max width=0.32\linewidth,max totalheight=4.8cm]{img/2009/shanghai_liberal/q16_fig1.png}

\begin{center}
\bitmapinclude[max width=0.88\linewidth,max totalheight=6.0cm]{img/2009/shanghai_liberal/q16_options.png}
\end{center}

\end{problem}

\begin{problem}
点 \(P(4, -2)\) 与圆 \(x^2 + y^2 = 4\) 上任一点连线的中点轨迹方程是（\quad）

\choices
    {\((x - 2)^2 + (y + 1)^2 = 1\)}
    {\((x - 2)^2 + (y + 1)^2 = 4\)}
    {\((x + 4)^2 + (y - 2)^2 = 4\)}
    {\((x + 2)^2 + (y - 1)^2 = 1\)}

\end{problem}

\begin{problem}
在发生某公共卫生事件期间，有专业机构认为该事件在一段时间没有发生规模群体感染的标志为“连续 10 天，每天新增疑似病例不超过 7 人”. 根据过去 10 天甲、乙、丙、丁四地新增疑似病例数据，一定符合该标志的是（\quad）

\choices
    {甲地：总体均值为 3，中位数为 4}
    {乙地：总体均值为 1，总体方差大于 0}
    {丙地：中位数为 2，众数为 3}
    {丁地：总体均值为 2，总体方差为 3}

\end{problem}

\setcounter{probnum}{19}
\section{解答题}

\begin{problem}
已知 \(\triangle ABC\) 的角 \(A\)、\(B\)、\(C\) 所对的边分别是 \(a\)、\(b\)、\(c\)，设向量 \(m = (a, b)\)，\(n = (\sin B, \sin A)\)，\(p = (b - 2, a - 2)\).

\begin{enumerate}[label=（\arabic*）]
\item 若 \(m \parallel n\)，求证：\(\triangle ABC\) 为等腰三角形；

\item 若 \(m \perp p\)，边长 \(c = 2\)，角 \(C = \frac{\pi}{3}\)，求 \(\triangle ABC\) 的面积.
\end{enumerate}

\end{problem}

\begin{problem}
有时可用函数 \(f(x) = \begin{cases} 0.1 + 15 \ln \frac{a}{a - x}, & x \leqslant 6, \\ \frac{x - 4.4}{x - 4}, & x > 6, \end{cases}\) 描述学习某学科知识的掌握程度，其中 \(x\) 表示某学科知识的学习次数（\(x \in \mathbf{N}^*\)），\(f(x)\) 表示对该学科知识的掌握程度，正实数 \(a\) 与学科知识有关.

\begin{enumerate}[label=（\arabic*）]
\item 证明：当 \(x \geqslant 7\) 时，掌握程度的增加量 \(f(x + 1) - f(x)\) 总是下降；

\item 根据经验，学科甲、乙、丙对应的 \(a\) 的取值范围分别为 \((115, 121]\)，\((121, 127]\)，\((121, 133]\). 当学习某学科知识 6 次时，掌握程度是 \(85\%\)，请确定相应的学科.
\end{enumerate}

\end{problem}

\begin{problem}
已知双曲线 \(C\) 的中心是原点，右焦点为 \(F(\sqrt{3}, 0)\)，一条渐近线 \(m: x + \sqrt{2}y = 0\)，设过点 \(A(-3\sqrt{2}, 0)\) 的直线 \(l\) 的方向向量 \(e = (1, k)\).

\begin{enumerate}[label=（\arabic*）]
\item 求双曲线 \(C\) 的方程；

\item 若过原点的直线 \(a \parallel l\)，且 \(a\) 与 \(l\) 的距离为 \(\sqrt{6}\)，求 \(k\) 的值；

\item 证明：当 \(k > \frac{\sqrt{2}}{2}\) 时，在双曲线 \(C\) 的右支上不存在点 \(Q\)，使之到直线 \(l\) 的距离为 \(\sqrt{6}\).
\end{enumerate}

\end{problem}

\begin{problem}
已知 \(\{a_n\}\) 是公差为 \(d\) 的等差数列，\(\{b_n\}\) 是公比为 \(q\) 的等比数列.

\begin{enumerate}[label=（\arabic*）]
\item 若 \(a_n = 3n + 1\)，是否存在 \(m, k \in \mathbf{N}^*\)，使 \(a_m + a_{m+1} = a_k\)？请说明理由；

\item 若 \(b_n = aq^n\)（\(a, q\) 为常数，且 \(aq \neq 0\)），对任意 \(m\) 存在 \(k\)，使 \(b_m \cdot b_{m+1} = b_k\)，试求 \(a, q\) 满足的充要条件；

\item 若 \(a_n = 2n + 1\)，\(b_n = 3^n\)，试确定所有的 \(p\)，使数列 \(\{b_n\}\) 中存在某个连续 \(p\) 项的和是数列 \(\{a_n\}\) 中的一项，请证明.
\end{enumerate}
\end{problem}
